\documentclass[a4paper,10pt,twoside]{article}
\usepackage{fontspec}
\setmainfont{Liberation Sans}
\setsansfont{Liberation Sans}
\renewcommand{\familydefault}{\sfdefault}
\usepackage{amsmath,amssymb}
\usepackage{geometry}
\usepackage{multicol}
\usepackage{graphicx}
\usepackage{fancyhdr}
\usepackage{lastpage}
\usepackage{booktabs}
\usepackage{listings}
\usepackage{eso-pic}
\usepackage{tikz}
\usepackage{array}
\usepackage{tabularx}
\usepackage{enumitem}
\usepackage{microtype}

\geometry{
    a4paper,
    top=1.8cm,
    bottom=1.8cm,
    left=1.2cm,
    right=1.2cm,
    headsep=14pt,
    footskip=20pt
}

\setlength{\columnsep}{16pt}
\setlength{\columnseprule}{0.35pt}

% Persistent Subtle Watermark on Every Page
\newsavebox{\watermarkbox}
\sbox{\watermarkbox}{%
    \begin{tabular}{c}
        \includegraphics[width=6.0cm]{/home/johan/GATE/GATE_CS_PYQ_DB/assets/studiva_butterfly_icon_transparent.jpeg} \\[-0.2em]
        {\fontsize{14}{17}\selectfont\bfseries Studiva}
    \end{tabular}%
}
\AddToShipoutPictureBG{%
    \begin{tikzpicture}[remember picture, overlay]
        \node[opacity=0.035] at (current page.center) {\usebox{\watermarkbox}};
    \end{tikzpicture}%
}

% Running Headers & Footers
\pagestyle{fancy}
\fancyhf{}
\fancyhead[C]{\small\textbf{Studiva GATE Question Bank}}
\fancyhead[LE]{\small\textit{Computer Science \& IT (2011--2025)}}
\fancyhead[RO]{\small\textit{Studiva Education}}
\fancyfoot[LE,RO]{\scriptsize\textbf{Studiva} $\cdot$ \texttt{www.studiva.co.in}}
\fancyfoot[C]{\small\textbf{Page \thepage\ of \pageref{LastPage}}}
\fancyfoot[RE,LO]{\scriptsize\texttt{contact@studiva.co.in}}
\renewcommand{\headrulewidth}{0.4pt}
\renewcommand{\footrulewidth}{0.4pt}

% Listings Style (Clean, professional)
\lstset{
    basicstyle=\scriptsize\ttfamily,
    numbers=none,
    frame=single,
    rulecolor=\color{black!30},
    tabsize=2,
    breaklines=true,
    showstringspaces=false
}

\begin{document}

% ==================== PAGE 1: FORMAL TITLE & COVER PAGE ====================
\thispagestyle{empty}
\vspace*{1.5cm}
\begin{center}
    \includegraphics[width=2.5cm]{/home/johan/GATE/GATE_CS_PYQ_DB/assets/studiva_butterfly_icon_transparent.jpeg}\\[1.5em]
    {\fontsize{24}{28}\selectfont\textbf{GATE COMPUTER SCIENCE \& IT}}\\[0.5em]
    {\fontsize{16}{20}\selectfont\textbf{PREVIOUS YEARS QUESTION BANK (2011 -- 2025)}}\\[0.6em]
    {\normalsize Detailed Examination Papers with Complete Official Answer Keys}\\[2.0em]
    \rule{0.65\textwidth}{0.8pt}\\[2.0em]
    {\large\textbf{Studiva GATE Examination Series}}\\[0.4em]
    {\small Comprehensive Authentic Question Bank for GATE Computer Science \& Information Technology}\\[2.8cm]

    \begin{tabular}{|ll|}
        \hline
        \textbf{Subject:} & Computer Science and Information Technology (CS/IT) \\
        \textbf{Period Covered:} & 2011 to 2025 \\
        \textbf{Examination Papers:} & 1 Complete Verified Paper Sets \\
        \textbf{Contents:} & Examination Questions \& Official Final Answer Keys \\
        \textbf{Published by:} & Studiva (\texttt{www.studiva.co.in}) \\
        \hline
    \end{tabular}

    \vfill
    {\small\textbf{Studiva $\cdot$ \texttt{www.studiva.co.in} $\cdot$ \texttt{contact@studiva.co.in}}}
\end{center}

\newpage

% ==================== PAGE 2: PUBLICATION IMPRINT & NOTES ====================
\thispagestyle{empty}
\vspace*{0.6cm}
\begin{center}
    {\large\textbf{PUBLICATION INFORMATION \& COPYRIGHT NOTICE}}\\[0.4em]
    \rule{0.6\textwidth}{0.5pt}
\end{center}

\vspace{1.2em}
\noindent\begin{tabular}{ll}
\textbf{Title:} & Studiva GATE Question Bank: Computer Science \& IT (2011--2025) \\
\textbf{Publisher:} & Studiva \\
\textbf{Website:} & \texttt{www.studiva.co.in} \\
\textbf{Official Inquiries:} & \texttt{contact@studiva.co.in} \\
\textbf{Edition:} & First Edition (2026) \\
\textbf{Syllabus Scope:} & Official GATE Syllabus (General Aptitude and CS / IT) \\
\end{tabular}

\vspace{1.8em}
\noindent\fbox{\parbox{\dimexpr\linewidth-2\fboxsep-2\fboxrule}{%
\textbf{COPYRIGHT NOTICE}\\[0.4em]
Copyright \copyright\ 2026 Studiva. All Rights Reserved.\\[0.4em]
No part of this publication may be reproduced, stored in a retrieval system, or transmitted in any form or by any means, electronic, mechanical, photocopying, recording, or otherwise, without prior written permission of Studiva, except for individual academic study.
}}

\vspace{1.5em}
\noindent\fbox{\parbox{\dimexpr\linewidth-2\fboxsep-2\fboxrule}{%
\textbf{PUBLISHER'S NOTE: ARCHIVAL QUESTION FORMAT}\\[0.4em]
This compendium has been meticulously curated and verified against official GATE question papers and final answer keys.
\par\medskip\noindent
\textbf{Notice on Image-Based Questions:}\\
Questions containing diagrams, circuit schematics, graphs, or image-based figures from the original examination papers have been omitted in this text edition due to non-readable format in historical source archives. All text-based conceptual, analytical, mathematical, algorithmic, and numerical questions are fully included with verified official keys.
}}

\vspace{1.5em}
\begin{center}
\textbf{EXAMINATION STRUCTURE \& PRACTICE GUIDELINES}
\end{center}
\begin{enumerate}[leftmargin=*]
    \item \textbf{Sectional Mark Division:} Each paper comprises General Aptitude (GA) and Computer Science (CS). Questions are clearly grouped into 1-mark and 2-mark blocks at the head of each section.
    \item \textbf{Question Types:} Includes Multiple Choice Questions (MCQ), Multiple Select Questions (MSQ), and Numerical Answer Type (NAT).
    \item \textbf{Official Answer Keys:} For self-evaluation without distraction, complete official answer keys are printed at the conclusion of each respective examination paper set.
\end{enumerate}

\newpage

% ==================== PAGE 3: TABLE OF CONTENTS ====================
\begin{center}
    {\Large\textbf{TABLE OF CONTENTS}}\\[0.4em]
    \rule{0.5\textwidth}{0.6pt}
\end{center}
\vspace{1em}

\noindent\begin{tabularx}{\linewidth}{|X|l|c|r|}
\hline
\textbf{Examination Paper} & \textbf{Session / Shift} & \textbf{Total Questions} & \textbf{Content Type} \\
\hline

GATE 2023 Computer Science & Shift CS & 54 & Questions + Answer Key \\ \hline
\end{tabularx}

\vspace{1.5em}
\noindent{\footnotesize\textit{Note: All questions in this edition have been typeset from verified digital records and cross-referenced with official answer keys released by the organising IIT/IISc institutes.}}

\newpage

% ==================== PAPER: 2023_CS ====================
\begin{center}
{\Large\textbf{GATE 2023 EXAMINATION}}\\[3pt]
{\normalsize\textbf{COMPUTER SCIENCE AND INFORMATION TECHNOLOGY}}\\[2pt]
{\small\textbf{Session / Shift: CS}}\\[5pt]
\rule{\linewidth}{0.8pt}\vspace{-4pt}
{\footnotesize\textbf{Duration: 180 Minutes} \hfill \textbf{Total Marks: 100} \hfill \textbf{Questions: 54}}\\
\vspace{-4pt}\rule{\linewidth}{0.4pt}
\end{center}
\vspace{4pt}
\noindent\fbox{\parbox{\dimexpr\linewidth-2\fboxsep-2\fboxrule}{%
\footnotesize\textbf{Examination Instructions:}\\[1pt]
1. Question Types: Multiple Choice Questions (MCQ), Multiple Select Questions (MSQ), and Numerical Answer Type (NAT).\\
2. Negative Marking Scheme: For 1-mark MCQ, $1/3$ mark is deducted; for 2-mark MCQ, $2/3$ mark is deducted. There is no negative marking for MSQ and NAT questions.\\
3. Mark Division: Questions are grouped by mark allocation within each section as specified below. Full official answer key is provided at the end of this paper.
}}
\vspace{10pt}
\begin{multicols}{2}

\vspace{8pt}
\noindent\textbf{\large\underline{GENERAL APTITUDE}}\\[2pt]
\noindent{\footnotesize\textit{[Section Block: 1-mark questions]}}
\vspace{4pt}

\noindent\textbf{Q.1} \hfill {\footnotesize\textsf{[MCQ]}}
\par\vspace{2pt}
{\bfseries We reached the station late, and \_\_\_\_\_\_\_ missed the train. \par}
\par\vspace{2pt}\noindent
\begin{tabular*}{\columnwidth}{@{\extracolsep{\fill}}p{0.46\columnwidth}p{0.46\columnwidth}@{}}
\textbf{(A)}~near & \textbf{(B)}~nearly \\
\addlinespace[2pt]
\textbf{(C)}~utterly & \textbf{(D)}~mostly
\end{tabular*}
\par\vspace{6pt}\hrule height 0.25pt\vspace{6pt}
\noindent\textbf{Q.2} \hfill {\footnotesize\textsf{[MCQ]}}
\par\vspace{2pt}
{\bfseries Kind : \_\_\_\_\_\_\_ : : Often : Frequently
(By word meaning) \par}
\par\vspace{2pt}\noindent
\begin{tabular*}{\columnwidth}{@{\extracolsep{\fill}}p{0.46\columnwidth}p{0.46\columnwidth}@{}}
\textbf{(A)}~Mean & \textbf{(B)}~Type \\
\addlinespace[2pt]
\textbf{(C)}~Cruel & \textbf{(D)}~Kindly
\end{tabular*}
\par\vspace{6pt}\hrule height 0.25pt\vspace{6pt}
\noindent\textbf{Q.3} \hfill {\footnotesize\textsf{[MCQ]}}
\par\vspace{2pt}
{\bfseries A series of natural numbers 𝐹1, 𝐹2, 𝐹3, 𝐹4, 𝐹5, 𝐹6, 𝐹7, … obeys 𝐹$n$+1 = 𝐹$n$+ 𝐹$n$−1
for all integers $n$$\ge$2 .
If 𝐹6 = 37, and 𝐹7 = 60, then what is 𝐹1 ? \par}
\par\vspace{2pt}\noindent
\begin{tabular*}{\columnwidth}{@{\extracolsep{\fill}}p{0.46\columnwidth}p{0.46\columnwidth}@{}}
\textbf{(A)}~4 & \textbf{(B)}~5 \\
\addlinespace[2pt]
\textbf{(C)}~8 & \textbf{(D)}~9
\end{tabular*}
\par\vspace{6pt}\hrule height 0.25pt\vspace{6pt}
\noindent\textbf{Q.4} \hfill {\footnotesize\textsf{[MCQ]}}
\par\vspace{2pt}
{\bfseries A survey for a certain year found that 90\% of pregnant women received medical
care at least once before giving birth. Of these women, 60\% received medical care
from doctors, while 40\% received medical care from other healthcare providers.
Given this information, which one of the following statements can be inferred with
certainty? \par}
\par\vspace{2pt}\noindent
\textbf{(A)}~More than half of the pregnant women received medical care at least once from a \newline doctor.\par\vspace{2pt}\noindent
\textbf{(B)}~Less than half of the pregnant women received medical care at least once from a \newline doctor.\par\vspace{2pt}\noindent
\textbf{(C)}~More than half of the pregnant women received medical care at most once from a \newline doctor.\par\vspace{2pt}\noindent
\textbf{(D)}~Less than half of the pregnant women received medical care at most once from a \newline doctor.
\par\vspace{6pt}\hrule height 0.25pt\vspace{6pt}
\noindent\textbf{Q.5} \hfill {\footnotesize\textsf{[MCQ]}}
\par\vspace{2pt}
{\bfseries Looking at the surface of a smooth 3-dimensional object from the outside, which
one of the following options is TRUE? \par}
\par\vspace{2pt}\noindent
\textbf{(A)}~The surface of the object must be concave everywhere.\par\vspace{2pt}\noindent
\textbf{(B)}~The surface of the object must be convex everywhere.\par\vspace{2pt}\noindent
\textbf{(C)}~The surface of the object may be concave in some places and convex in other places.\par\vspace{2pt}\noindent
\textbf{(D)}~The object can have edges, but no corners.
\par\vspace{6pt}\hrule height 0.25pt\vspace{6pt}

\vspace{6pt}
\noindent{\footnotesize\textit{[Section Block: 2-mark questions]}}
\vspace{3pt}

\noindent\textbf{Q.6} \hfill {\footnotesize\textsf{[MCQ]}}
\par\vspace{2pt}
{\bfseries The country of Zombieland is in distress since more than 75\% of its working
population is suffering from serious health issues. Studies conducted by competent
health experts concluded that a complete lack of physical exercise among its
working population was one of the leading causes of their health issues. As one of
the measures to address the problem, the Government of Zombieland has decided
to provide monetary incentives to those who ride bicycles to work.
\par\smallskip
Based only on the information provided above, which one of the following
statements can be logically inferred with certainty? \par}
\par\vspace{2pt}\noindent
\textbf{(A)}~All the working population of Zombieland will henceforth ride bicycles to work.\par\vspace{2pt}\noindent
\textbf{(B)}~Riding bicycles will ensure that all of the working population of Zombieland is free \newline of health issues.\par\vspace{2pt}\noindent
\textbf{(C)}~The health experts suggested to the Government of Zombieland to declare riding \newline bicycles as mandatory.\par\vspace{2pt}\noindent
\textbf{(D)}~The Government of Zombieland believes that riding bicycles is a form of physical \newline exercise.
\par\vspace{6pt}\hrule height 0.25pt\vspace{6pt}
\noindent\textbf{Q.7} \hfill {\footnotesize\textsf{[MCQ]}}
\par\vspace{2pt}
{\bfseries Consider two functions of time (𝑡),
𝑓(𝑡) = 0.01 𝑡2
𝑔(𝑡) = 4 𝑡
where 0 < 𝑡< $\infty$.
Now consider the following two statements:
(i)
For some 𝑡> 0, 𝑔(𝑡) > 𝑓(𝑡).
(ii)
There exists a 𝑇, such that 𝑓(𝑡) > 𝑔(𝑡) for all 𝑡> 𝑇.
Which one of the following options is TRUE? \par}
\par\vspace{2pt}\noindent
\textbf{(A)}~only (i) is correct\par\vspace{2pt}\noindent
\textbf{(B)}~only (ii) is correct\par\vspace{2pt}\noindent
\textbf{(C)}~both (i) and (ii) are correct\par\vspace{2pt}\noindent
\textbf{(D)}~neither (i) nor (ii) is correct
\par\vspace{6pt}\hrule height 0.25pt\vspace{6pt}
\noindent\textbf{Q.8} \hfill {\footnotesize\textsf{[MCQ]}}
\par\vspace{2pt}
{\bfseries Which one of the following sentence sequences creates a coherent narrative?
(i)
Once on the terrace, on her way to her small room in the corner, she
notices the man right away.
(ii)
She begins to pant by the time she has climbed all the stairs.
(iii)
Mina has bought vegetables and rice at the market, so her bags are
heavy.
(iv)
He was leaning against the parapet, watching the traffic below. \par}
\par\vspace{2pt}\noindent
\textbf{(A)}~(i), (ii), (iv), (iii)\par\vspace{2pt}\noindent
\textbf{(B)}~(ii), (iii), (i), (iv)\par\vspace{2pt}\noindent
\textbf{(C)}~(iv), (ii), (i), (iii)\par\vspace{2pt}\noindent
\textbf{(D)}~(iii), (ii), (i), (iv)
\par\vspace{6pt}\hrule height 0.25pt\vspace{6pt}
\noindent\textbf{Q.9} \hfill {\footnotesize\textsf{[MCQ]}}
\par\vspace{2pt}
{\bfseries ( )
f x and ( )
g y are functions of x and y, respectively, and
( )
( )
f x
g y
=
for all real
values of $x$ and $y$. Which one of the following options is necessarily TRUE for all
x and y? \par}
\par\vspace{2pt}\noindent
\textbf{(A)}~( ) \newline 0 \newline f x = \newline and ( ) \newline 0 \newline g y =\par\vspace{2pt}\noindent
\textbf{(B)}~( ) \newline ( ) \newline constant \newline f x \newline g y \newline = \newline =\par\vspace{2pt}\noindent
\textbf{(C)}~( ) \newline constant \newline f x $\ne$ \newline and ( ) \newline constant \newline g y $\ne$\par\vspace{2pt}\noindent
\textbf{(D)}~( ) \newline ( ) \newline ( ) \newline ( ) \newline = \newline  \newline f x \newline g y \newline f x \newline g y \newline + \newline −
\par\vspace{6pt}\hrule height 0.25pt\vspace{6pt}
\noindent\textbf{Q.10} \hfill {\footnotesize\textsf{[MCQ]}}
\par\vspace{2pt}
{\bfseries Which one of the options best describes the transformation of the 2-dimensional
figure P to Q, and then to R, as shown? \par}
\par\vspace{2pt}\noindent
\textbf{(A)}~Only S1 is TRUE.\par\vspace{2pt}\noindent
\textbf{(B)}~Only S1 and S2 are TRUE.\par\vspace{2pt}\noindent
\textbf{(C)}~S1, S2, and S3 are all TRUE.\par\vspace{2pt}\noindent
\textbf{(D)}~Only S1 and S3 are TRUE. \newline CS
\par\vspace{6pt}\hrule height 0.25pt\vspace{6pt}

\vspace{8pt}
\noindent\textbf{\large\underline{COMPUTER SCIENCE AND INFORMATION TECHNOLOGY}}\\[2pt]
\noindent{\footnotesize\textit{[Section Block: 1-mark questions]}}
\vspace{4pt}

\noindent\textbf{Q.12} \hfill {\footnotesize\textsf{[MCQ]}}
\par\vspace{2pt}
{\bfseries Which one of the following sequences when stored in an array at locations
A[1], . . . , A[10] forms a max-heap? \par}
\par\vspace{2pt}\noindent
\textbf{(A)}~23, 17, 10, 6, 13, 14, 1, 5, 7, 12\par\vspace{2pt}\noindent
\textbf{(B)}~23, 17, 14, 7, 13, 10, 1, 5, 6, 12\par\vspace{2pt}\noindent
\textbf{(C)}~23, 17, 14, 6, 13, 10, 1, 5, 7, 15\par\vspace{2pt}\noindent
\textbf{(D)}~23, 14, 17, 1, 10, 13, 16, 12, 7, 5 \newline CS
\par\vspace{6pt}\hrule height 0.25pt\vspace{6pt}
\noindent\textbf{Q.13} \hfill {\footnotesize\textsf{[MCQ]}}
\par\vspace{2pt}
{\bfseries Let SLLdel be a function that deletes a node in a singly-linked list given a pointer
to the node and a pointer to the head of the list. Similarly, let DLLdel be another
function that deletes a node in a doubly-linked list given a pointer to the node and
a pointer to the head of the list.
Let n denote the number of nodes in each of the linked lists. Which one of the
following choices is TRUE about the worst-case time complexity of SLLdel and
DLLdel? \par}
\par\vspace{2pt}\noindent
\textbf{(A)}~SLLdel is O(1) and DLLdel is O(n)\par\vspace{2pt}\noindent
\textbf{(B)}~Both SLLdel and DLLdel are O(log(n))\par\vspace{2pt}\noindent
\textbf{(C)}~Both SLLdel and DLLdel are O(1)\par\vspace{2pt}\noindent
\textbf{(D)}~SLLdel is O(n) and DLLdel is O(1) \newline CS
\par\vspace{6pt}\hrule height 0.25pt\vspace{6pt}
\noindent\textbf{Q.15} \hfill {\footnotesize\textsf{[MCQ]}}
\par\vspace{2pt}
{\bfseries The Lucas sequence Ln is defined by the recurrence relation:
Ln = Ln−1 + Ln−2,
for
n $\ge$3,
with L1 = 1 and L2 = 3.
Which one of the options given is TRUE? \par}
\par\vspace{2pt}\noindent
\textbf{(A)}~Ln = \newline  \newline 1 + \newline $\sqrt{\vphantom{x}}$ \newline 5 \newline 2 \newline n \newline + \newline  \newline 1 − \newline $\sqrt{\vphantom{x}}$ \newline 5 \newline 2 \newline n\par\vspace{2pt}\noindent
\textbf{(B)}~Ln = \newline  \newline 1 + \newline $\sqrt{\vphantom{x}}$ \newline 5 \newline 2 \newline n \newline − \newline  \newline 1 − \newline $\sqrt{\vphantom{x}}$ \newline 5 \newline 3 \newline n\par\vspace{2pt}\noindent
\textbf{(C)}~Ln = \newline  \newline 1 + \newline $\sqrt{\vphantom{x}}$ \newline 5 \newline 2 \newline n \newline + \newline  \newline 1 − \newline $\sqrt{\vphantom{x}}$ \newline 5 \newline 3 \newline n\par\vspace{2pt}\noindent
\textbf{(D)}~Ln = \newline  \newline 1 + \newline $\sqrt{\vphantom{x}}$ \newline 5 \newline 2 \newline n \newline − \newline  \newline 1 − \newline $\sqrt{\vphantom{x}}$ \newline 5 \newline 2 \newline n \newline CS
\par\vspace{6pt}\hrule height 0.25pt\vspace{6pt}
\noindent\textbf{Q.16} \hfill {\footnotesize\textsf{[MCQ]}}
\par\vspace{2pt}
{\bfseries Which one of the options given below refers to the degree (or arity) of a relation in
relational database systems? \par}
\par\vspace{2pt}\noindent
\textbf{(A)}~Number of attributes of its relation schema.\par\vspace{2pt}\noindent
\textbf{(B)}~Number of tuples stored in the relation.\par\vspace{2pt}\noindent
\textbf{(C)}~Number of entries in the relation.\par\vspace{2pt}\noindent
\textbf{(D)}~Number of distinct domains of its relation schema. \newline CS
\par\vspace{6pt}\hrule height 0.25pt\vspace{6pt}
\noindent\textbf{Q.17} \hfill {\footnotesize\textsf{[MCQ]}}
\par\vspace{2pt}
{\bfseries Suppose two hosts are connected by a point-to-point link and they are configured
to use Stop-and-Wait protocol for reliable data transfer. Identify in which one of
the following scenarios, the utilization of the link is the lowest. \par}
\par\vspace{2pt}\noindent
\textbf{(A)}~Longer link length and lower transmission rate\par\vspace{2pt}\noindent
\textbf{(B)}~Longer link length and higher transmission rate\par\vspace{2pt}\noindent
\textbf{(C)}~Shorter link length and lower transmission rate\par\vspace{2pt}\noindent
\textbf{(D)}~Shorter link length and higher transmission rate \newline CS
\par\vspace{6pt}\hrule height 0.25pt\vspace{6pt}
\noindent\textbf{Q.18} \hfill {\footnotesize\textsf{[MCQ]}}
\par\vspace{2pt}
{\bfseries Let
A =
⎡
⎢⎢⎣
1
2
3
4
4
1
2
3
3
4
1
2
2
3
4
1
⎤
⎥⎥⎦
and
B =
⎡
⎢⎢⎣
3
4
1
2
4
1
2
3
1
2
3
4
2
3
4
1
⎤
⎥⎥⎦.
Let det(A) and det(B) denote the determinants of the matrices A and B,
respectively.
Which one of the options given below is TRUE? \par}
\par\vspace{2pt}\noindent
\textbf{(A)}~det(A) = det(B)\par\vspace{2pt}\noindent
\textbf{(B)}~det(B) = −det(A)\par\vspace{2pt}\noindent
\textbf{(C)}~det(A) = 0\par\vspace{2pt}\noindent
\textbf{(D)}~det(AB) = det(A) + det(B) \newline CS
\par\vspace{6pt}\hrule height 0.25pt\vspace{6pt}
\noindent\textbf{Q.20} \hfill {\footnotesize\textsf{[MCQ]}}
\par\vspace{2pt}
{\bfseries CS
\par\smallskip
An algorithm has to store several keys generated by an adversary in a hash
table. The adversary is malicious who tries to maximize the number of collisions.
Let k be the number of keys, m be the number of slots in the hash table, and k > m.
Which one of the following is the best hashing strategy to counteract the adversary? \par}
\par\vspace{2pt}\noindent
\textbf{(A)}~Division method, i.e., use the hash function h(k) = k mod m.\par\vspace{2pt}\noindent
\textbf{(B)}~Multiplication method, i.e., use the hash function h(k) = ⌊m(kA −⌊kA⌋)⌋, where \newline A is a carefully chosen constant.\par\vspace{2pt}\noindent
\textbf{(C)}~Universal hashing method.\par\vspace{2pt}\noindent
\textbf{(D)}~If k is a prime number, use Division method. Otherwise, use Multiplication method. \newline CS
\par\vspace{6pt}\hrule height 0.25pt\vspace{6pt}
\noindent\textbf{Q.22} \hfill {\footnotesize\textsf{[MSQ]}}
\par\vspace{2pt}
{\bfseries Which one or more of the following need to be saved on a context switch from one
thread (T1) of a process to another thread (T2) of the same process? \par}
\par\vspace{2pt}\noindent
\textbf{(A)}~Page table base register\par\vspace{2pt}\noindent
\textbf{(B)}~Stack pointer\par\vspace{2pt}\noindent
\textbf{(C)}~Program counter\par\vspace{2pt}\noindent
\textbf{(D)}~General purpose registers \newline CS
\par\vspace{6pt}\hrule height 0.25pt\vspace{6pt}
\noindent\textbf{Q.23} \hfill {\footnotesize\textsf{[MSQ]}}
\par\vspace{2pt}
{\bfseries Which one or more of the following options guarantee that a computer system will
transition from user mode to kernel mode? \par}
\par\vspace{2pt}\noindent
\textbf{(A)}~Function Call\par\vspace{2pt}\noindent
\textbf{(B)}~malloc Call\par\vspace{2pt}\noindent
\textbf{(C)}~Page Fault\par\vspace{2pt}\noindent
\textbf{(D)}~System Call \newline CS
\par\vspace{6pt}\hrule height 0.25pt\vspace{6pt}
\noindent\textbf{Q.24} \hfill {\footnotesize\textsf{[MSQ]}}
\par\vspace{2pt}
{\bfseries Which of the following statements is/are CORRECT? \par}
\par\vspace{2pt}\noindent
\textbf{(A)}~The intersection of two regular languages is regular.\par\vspace{2pt}\noindent
\textbf{(B)}~The intersection of two context-free languages is context-free.\par\vspace{2pt}\noindent
\textbf{(C)}~The intersection of two recursive languages is recursive.\par\vspace{2pt}\noindent
\textbf{(D)}~The intersection of two recursively enumerable languages is recursively enumerable. \newline CS
\par\vspace{6pt}\hrule height 0.25pt\vspace{6pt}
\noindent\textbf{Q.25} \hfill {\footnotesize\textsf{[MSQ]}}
\par\vspace{2pt}
{\bfseries Which of the following statements is/are INCORRECT about
the OSPF (Open Shortest Path First) routing protocol used in the Internet? \par}
\par\vspace{2pt}\noindent
\textbf{(A)}~OSPF implements Bellman-Ford algorithm to find shortest paths.\par\vspace{2pt}\noindent
\textbf{(B)}~OSPF uses Dijkstra's shortest path algorithm to implement least-cost path routing.\par\vspace{2pt}\noindent
\textbf{(C)}~OSPF is used as an inter-domain routing protocol.\par\vspace{2pt}\noindent
\textbf{(D)}~OSPF implements hierarchical routing. \newline CS
\par\vspace{6pt}\hrule height 0.25pt\vspace{6pt}
\noindent\textbf{Q.26} \hfill {\footnotesize\textsf{[MSQ]}}
\par\vspace{2pt}
{\bfseries Geetha has a conjecture about integers, which is of the form
∀x
P(x) =$\implies$∃yQ(x, y)
,
where P is a statement about integers, and Q is a statement about pairs of integers.
Which of the following (one or more) option(s) would imply Geetha's conjecture? \par}
\par\vspace{2pt}\noindent
\textbf{(A)}~∃x \newline P(x) ∧∀yQ(x, y)\par\vspace{2pt}\noindent
\textbf{(B)}~∀x∀yQ(x, y)\par\vspace{2pt}\noindent
\textbf{(C)}~∃y∀x \newline P(x) =$\implies$Q(x, y)\par\vspace{2pt}\noindent
\textbf{(D)}~∃x \newline P(x) ∧∃yQ(x, y) \newline CS
\par\vspace{6pt}\hrule height 0.25pt\vspace{6pt}
\noindent\textbf{Q.27} \hfill {\footnotesize\textsf{[MSQ]}}
\par\vspace{2pt}
{\bfseries Which one or more of the following CPU scheduling algorithms can potentially
cause starvation? \par}
\par\vspace{2pt}\noindent
\textbf{(A)}~First-in First-Out\par\vspace{2pt}\noindent
\textbf{(B)}~Round Robin\par\vspace{2pt}\noindent
\textbf{(C)}~Priority Scheduling\par\vspace{2pt}\noindent
\textbf{(D)}~Shortest Job First \newline CS
\par\vspace{6pt}\hrule height 0.25pt\vspace{6pt}
\noindent\textbf{Q.28} \hfill {\footnotesize\textsf{[MSQ]}}
\par\vspace{2pt}
{\bfseries Let
f(x) = x3 + 15x2 −33x −36
be a real-valued function.
Which of the following statements is/are TRUE? \par}
\par\vspace{2pt}\noindent
\textbf{(A)}~f(x) does not have a local maximum.\par\vspace{2pt}\noindent
\textbf{(B)}~f(x) has a local maximum.\par\vspace{2pt}\noindent
\textbf{(C)}~f(x) does not have a local minimum.\par\vspace{2pt}\noindent
\textbf{(D)}~f(x) has a local minimum. \newline CS
\par\vspace{6pt}\hrule height 0.25pt\vspace{6pt}
\noindent\textbf{Q.29} \hfill {\footnotesize\textsf{[MSQ]}}
\par\vspace{2pt}
{\bfseries Let f and g be functions of natural numbers given by f(n) = n and g(n) = n2.
Which of the following statements is/are TRUE? \par}
\par\vspace{2pt}\noindent
\textbf{(A)}~f $\in$O(g)\par\vspace{2pt}\noindent
\textbf{(B)}~f $\in$$\Omega$(g)\par\vspace{2pt}\noindent
\textbf{(C)}~f $\in$o(g)\par\vspace{2pt}\noindent
\textbf{(D)}~f $\in$Θ(g) \newline CS
\par\vspace{6pt}\hrule height 0.25pt\vspace{6pt}
\noindent\textbf{Q.31} \hfill {\footnotesize\textsf{[NAT]}}
\par\vspace{2pt}
{\bfseries The value of the definite integral
3
−3
2
−2
1
−1
4x2y −z3
dz dy dx
is
. (Rounded offto the nearest integer)
CS \par}
\par\vspace{3pt}\noindent\textbf{Answer:} \rule{3.2cm}{0.4pt}
\par\vspace{6pt}\hrule height 0.25pt\vspace{6pt}
\noindent\textbf{Q.32} \hfill {\footnotesize\textsf{[NAT]}}
\par\vspace{2pt}
{\bfseries A particular number is written as 132 in radix-4 representation. The same number
in radix-5 representation is
.
CS \par}
\par\vspace{3pt}\noindent\textbf{Answer:} \rule{3.2cm}{0.4pt}
\par\vspace{6pt}\hrule height 0.25pt\vspace{6pt}
\noindent\textbf{Q.33} \hfill {\footnotesize\textsf{[NAT]}}
\par\vspace{2pt}
{\bfseries Consider a 3-stage pipelined processor having a delay of 10 ns (nanoseconds),
20 ns, and 14 ns, for the first, second, and the third stages, respectively. Assume
that there is no other delay and the processor does not suffer from any pipeline
hazards. Also assume that one instruction is fetched every cycle.
The total execution time for executing 100 instructions on this processor is
ns.
CS \par}
\par\vspace{3pt}\noindent\textbf{Answer:} \rule{3.2cm}{0.4pt}
\par\vspace{6pt}\hrule height 0.25pt\vspace{6pt}
\noindent\textbf{Q.34} \hfill {\footnotesize\textsf{[NAT]}}
\par\vspace{2pt}
{\bfseries A keyboard connected to a computer is used at a rate of 1 keystroke per second.
The computer system polls the keyboard every 10 ms (milli seconds) to check for
a keystroke and consumes 100 $\mu$s (micro seconds) for each poll. If it is determined
after polling that a key has been pressed, the system consumes an additional 200
$\mu$s to process the keystroke. Let T1 denote the fraction of a second spent in polling
and processing a keystroke.
In an alternative implementation, the system uses interrupts instead of polling.
An interrupt is raised for every keystroke. It takes a total of 1 ms for servicing an
interrupt and processing a keystroke. Let T2 denote the fraction of a second spent
in servicing the interrupt and processing a keystroke.
The ratio T1
T2
is
. (Rounded offto one decimal place)
CS \par}
\par\vspace{3pt}\noindent\textbf{Answer:} \rule{3.2cm}{0.4pt}
\par\vspace{6pt}\hrule height 0.25pt\vspace{6pt}
\noindent\textbf{Q.35} \hfill {\footnotesize\textsf{[NAT]}}
\par\vspace{2pt}
{\bfseries The integer value printed by the ANSI-C program given below is
.
\#include<stdio.h>
int funcp()\{
static int x = 1;
x++;
return x;
\}
int main()\{
int x,y;
x = funcp();
y = funcp()+x;
printf("\%d\n", (x+y));
return 0;
\}
CS \par}
\par\vspace{3pt}\noindent\textbf{Answer:} \rule{3.2cm}{0.4pt}
\par\vspace{6pt}\hrule height 0.25pt\vspace{6pt}
\noindent\textbf{Q.36} \hfill {\footnotesize\textsf{[MCQ]}}
\par\vspace{2pt}
{\bfseries CS
\par\smallskip
Consider the following program:
int main()
\{
f1();
f2(2);
f3();
return(0);
\}
int f1()
\{
return(1);
\}
int f2(int X)
\{
f3();
if (X==1)
return f1();
else
return (X*f2(X-1));
\}
int f3()
\{
return(5);
\}
Which one of the following options represents the activation tree corresponding to
the main function? \par}
\par\vspace{2pt}\noindent
\textbf{(A)}~main \newline f1 f2 f3 \newline f3 \newline f2 \newline f3 \newline f1\par\vspace{2pt}\noindent
\textbf{(B)}~main \newline f2 f3 \newline f1 \newline f1 \newline f3\par\vspace{2pt}\noindent
\textbf{(C)}~main \newline f2 \newline f1 \newline f3 \newline f1\par\vspace{2pt}\noindent
\textbf{(D)}~main \newline f1 f2 f3 \newline f3 \newline f2 \newline f1 \newline CS
\par\vspace{6pt}\hrule height 0.25pt\vspace{6pt}

\vspace{6pt}
\noindent{\footnotesize\textit{[Section Block: 2-mark questions]}}
\vspace{3pt}

\noindent\textbf{Q.38} \hfill {\footnotesize\textsf{[MCQ]}}
\par\vspace{2pt}
{\bfseries CS
\par\smallskip
Consider the two functions incr and decr shown below.
incr()\{
wait(s);
X = X+1;
signal(s);
\}
decr()\{
wait(s);
X = X-1;
signal(s);
\}
There are 5 threads each invoking incr once, and 3 threads each invoking decr
once, on the same shared variable X. The initial value of X is 10.
Suppose there are two implementations of the semaphore s, as follows:
I-1: s is a binary semaphore initialized to 1.
I-2: s is a counting semaphore initialized to 2.
Let V1, V2 be the values of X at the end of execution of all the threads with
implementations I-1, I-2, respectively.
Which one of the following choices corresponds to the minimum possible values of
V1, V2, respectively? \par}
\par\vspace{2pt}\noindent
\textbf{(A)}~15, 7\par\vspace{2pt}\noindent
\textbf{(B)}~7, 7\par\vspace{2pt}\noindent
\textbf{(C)}~12, 7\par\vspace{2pt}\noindent
\textbf{(D)}~12, 8 \newline CS
\par\vspace{6pt}\hrule height 0.25pt\vspace{6pt}
\noindent\textbf{Q.39} \hfill {\footnotesize\textsf{[MCQ]}}
\par\vspace{2pt}
{\bfseries Consider the context-free grammar G below
S
$\to$
aSb | X
X
$\to$
aX | Xb | a | b ,
where S and X are non-terminals, and a and b are terminal symbols. The starting
non-terminal is S.
Which one of the following statements is CORRECT? \par}
\par\vspace{2pt}\noindent
\textbf{(A)}~The language generated by G is (a + b)∗\par\vspace{2pt}\noindent
\textbf{(B)}~The language generated by G is a∗(a + b)b∗\par\vspace{2pt}\noindent
\textbf{(C)}~The language generated by G is a∗b∗(a + b)\par\vspace{2pt}\noindent
\textbf{(D)}~The language generated by G is not a regular language \newline CS
\par\vspace{6pt}\hrule height 0.25pt\vspace{6pt}
\noindent\textbf{Q.41} \hfill {\footnotesize\textsf{[MCQ]}}
\par\vspace{2pt}
{\bfseries Consider the given C-code and its corresponding assembly code, with a few
operands U1--U4 being unknown. Some useful information as well as the semantics
of each unique assembly instruction is annotated as inline comments in the code.
The memory is byte-addressable.
//C-code
int a[10], b[10], i;
// int is 32-bit
for (i=0; i<10;i++)
a[i] = b[i] * 8;
;assembly-code (; indicates comments)
;r1-r5 are 32-bit integer registers
;initialize r1=0, r2=10
;initialize r3, r4 with base address of a, b
L01: jeq r1, r2, end
;if(r1==r2) goto end
L02: lw
r5, 0(r4)
;r5 <- Memory[r4+0]
L03: shl r5, r5, U1
;r5 <- r5 << U1
L04: sw
r5, 0(r3)
;Memory[r3+0] <- r5
L05: add r3, r3, U2
;r3 <- r3+U2
L06: add r4, r4, U3
L07: add r1, r1, 1
L08: jmp U4
;goto U4
L09: end
Which one of the following options is a CORRECT replacement for operands in
the position (U1, U2, U3, U4) in the above assembly code? \par}
\par\vspace{2pt}\noindent
\textbf{(A)}~(8, 4, 1, L02)\par\vspace{2pt}\noindent
\textbf{(B)}~(3, 4, 4, L01)\par\vspace{2pt}\noindent
\textbf{(C)}~(8, 1, 1, L02)\par\vspace{2pt}\noindent
\textbf{(D)}~(3, 1, 1, L01) \newline CS
\par\vspace{6pt}\hrule height 0.25pt\vspace{6pt}
\noindent\textbf{Q.45} \hfill {\footnotesize\textsf{[MCQ]}}
\par\vspace{2pt}
{\bfseries CS
\par\smallskip
Consider the IEEE-754 single precision floating point numbers
P=0xC1800000 and Q=0x3F5C2EF4.
Which one of the following corresponds to the product of these numbers
(i.e., P $\times$ Q), represented in the IEEE-754 single precision format? \par}
\par\vspace{2pt}\noindent
\textbf{(A)}~0x404C2EF4\par\vspace{2pt}\noindent
\textbf{(B)}~0x405C2EF4\par\vspace{2pt}\noindent
\textbf{(C)}~0xC15C2EF4\par\vspace{2pt}\noindent
\textbf{(D)}~0xC14C2EF4 \newline CS
\par\vspace{6pt}\hrule height 0.25pt\vspace{6pt}
\noindent\textbf{Q.46} \hfill {\footnotesize\textsf{[MCQ]}}
\par\vspace{2pt}
{\bfseries Let A be a priority queue for maintaining a set of elements.
Suppose A is
implemented using a max-heap data structure. The operation Extract-Max(A)
extracts and deletes the maximum element from A. The operation Insert(A,key)
inserts a new element key in A. The properties of a max-heap are preserved at the
end of each of these operations.
When A contains n elements, which one of the following statements about the worst
case running time of these two operations is TRUE? \par}
\par\vspace{2pt}\noindent
\textbf{(A)}~Both Extract-Max(A) and Insert(A,key) run in O(1).\par\vspace{2pt}\noindent
\textbf{(B)}~Both Extract-Max(A) and Insert(A,key) run in O(log(n)).\par\vspace{2pt}\noindent
\textbf{(C)}~Extract-Max(A) runs in O(1) whereas Insert(A,key) runs in O(n).\par\vspace{2pt}\noindent
\textbf{(D)}~Extract-Max(A) runs in O(1) whereas Insert(A,key) runs in O(log(n)). \newline CS
\par\vspace{6pt}\hrule height 0.25pt\vspace{6pt}
\noindent\textbf{Q.47} \hfill {\footnotesize\textsf{[MCQ]}}
\par\vspace{2pt}
{\bfseries Consider the C function foo and the binary tree shown.
typedef struct node \{
int val;
struct node *left, *right;
\} node;
int foo(node *p) \{
int retval;
if (p == NULL)
return 0;
else \{
retval = p->val + foo(p->left) + foo(p->right);
printf("\%d ", retval);
return retval;
\}
\}
10
5
3
8
11
13
When foo is called with a pointer to the root node of the given binary tree, what
will it print? \par}
\par\vspace{2pt}\noindent
\textbf{(A)}~3 8 5 13 11 10\par\vspace{2pt}\noindent
\textbf{(B)}~3 5 8 10 11 13\par\vspace{2pt}\noindent
\textbf{(C)}~3 8 16 13 24 50\par\vspace{2pt}\noindent
\textbf{(D)}~3 16 8 50 24 13 \newline CS
\par\vspace{6pt}\hrule height 0.25pt\vspace{6pt}
\noindent\textbf{Q.48} \hfill {\footnotesize\textsf{[MCQ]}}
\par\vspace{2pt}
{\bfseries Let U = \{1, 2, . . . , n\}, where n is a large positive integer greater than 1000. Let
k be a positive integer less than n. Let A, B be subsets of U with |A| = |B| = k
and A $\cap$B = ∅. We say that a permutation of U separates A from B if one of the
following is true.
- All members of A appear in the permutation before any of the members of B.
- All members of B appear in the permutation before any of the members of A.
How many permutations of U separate A from B? \par}
\par\vspace{2pt}\noindent
\textbf{(A)}~n!\par\vspace{2pt}\noindent
\textbf{(B)}~ n \newline 2k \newline  \newline (n −2k)!\par\vspace{2pt}\noindent
\textbf{(C)}~ n \newline 2k \newline  \newline (n −2k)! (k!)2\par\vspace{2pt}\noindent
\textbf{(D)}~2 \newline  n \newline 2k \newline  \newline (n −2k)! (k!)2 \newline CS
\par\vspace{6pt}\hrule height 0.25pt\vspace{6pt}
\noindent\textbf{Q.49} \hfill {\footnotesize\textsf{[MSQ]}}
\par\vspace{2pt}
{\bfseries Let f : A $\to$B be an onto (or surjective) function, where A and B are nonempty
sets. Define an equivalence relation ∼on the set A as
a1 ∼a2 if f(a1) = f(a2) ,
where a1, a2 $\in$A. Let E = \{[x] : x $\in$A\} be the set of all the equivalence classes
under ∼. Define a new mapping F : E $\to$B as
F([x]) = f(x),
for all the equivalence classes [x] in E.
Which of the following statements is/are TRUE? \par}
\par\vspace{2pt}\noindent
\textbf{(A)}~F is NOT well-defined.\par\vspace{2pt}\noindent
\textbf{(B)}~F is an onto (or surjective) function.\par\vspace{2pt}\noindent
\textbf{(C)}~F is a one-to-one (or injective) function.\par\vspace{2pt}\noindent
\textbf{(D)}~F is a bijective function. \newline CS
\par\vspace{6pt}\hrule height 0.25pt\vspace{6pt}
\noindent\textbf{Q.50} \hfill {\footnotesize\textsf{[MSQ]}}
\par\vspace{2pt}
{\bfseries Suppose you are asked to design a new reliable byte-stream transport protocol
like TCP. This protocol, named myTCP, runs over a 100 Mbps network with Round
Trip Time of 150 milliseconds and the maximum segment lifetime of 2 minutes.
Which of the following is/are valid lengths of the Sequence Number field in the
myTCP header? \par}
\par\vspace{2pt}\noindent
\textbf{(A)}~30 bits\par\vspace{2pt}\noindent
\textbf{(B)}~32 bits\par\vspace{2pt}\noindent
\textbf{(C)}~34 bits\par\vspace{2pt}\noindent
\textbf{(D)}~36 bits \newline CS
\par\vspace{6pt}\hrule height 0.25pt\vspace{6pt}
\noindent\textbf{Q.51} \hfill {\footnotesize\textsf{[MSQ]}}
\par\vspace{2pt}
{\bfseries Let X be a set and 2X denote the powerset of X.
Define a binary operation Δ on 2X as follows:
AΔB = (A −B) $\cup$(B −A) .
Let H = (2X, Δ). Which of the following statements about H is/are correct? \par}
\par\vspace{2pt}\noindent
\textbf{(A)}~H is a group.\par\vspace{2pt}\noindent
\textbf{(B)}~Every element in H has an inverse, but H is NOT a group.\par\vspace{2pt}\noindent
\textbf{(C)}~For every A $\in$2X, the inverse of A is the complement of A.\par\vspace{2pt}\noindent
\textbf{(D)}~For every A $\in$2X, the inverse of A is A. \newline CS
\par\vspace{6pt}\hrule height 0.25pt\vspace{6pt}
\noindent\textbf{Q.52} \hfill {\footnotesize\textsf{[MSQ]}}
\par\vspace{2pt}
{\bfseries Suppose in a web browser, you click on the www.gate-2023.in URL. The browser
cache is empty. The IP address for this URL is not cached in your local host, so
a DNS lookup is triggered (by the local DNS server deployed on your local host)
over the 3-tier DNS hierarchy in an iterative mode. No resource records are cached
anywhere across all DNS servers.
Let RTT denote the round trip time between your local host and DNS servers
in the DNS hierarchy. The round trip time between the local host and the web
server hosting www.gate-2023.in is also equal to RTT. The HTML file associated
with the URL is small enough to have negligible transmission time and negligible
rendering time by your web browser, which references 10 equally small objects on
the same web server.
Which of the following statements is/are CORRECT about the minimum elapsed
time between clicking on the URL and your browser fully rendering it? \par}
\par\vspace{2pt}\noindent
\textbf{(A)}~7 RTTs, in case of non-persistent HTTP with 5 parallel TCP connections.\par\vspace{2pt}\noindent
\textbf{(B)}~5 RTTs, in case of persistent HTTP with pipelining.\par\vspace{2pt}\noindent
\textbf{(C)}~9 RTTs, in case of non-persistent HTTP with 5 parallel TCP connections.\par\vspace{2pt}\noindent
\textbf{(D)}~6 RTTs, in case of persistent HTTP with pipelining. \newline CS
\par\vspace{6pt}\hrule height 0.25pt\vspace{6pt}
\noindent\textbf{Q.53} \hfill {\footnotesize\textsf{[MSQ]}}
\par\vspace{2pt}
{\bfseries Consider a random experiment where two fair coins are tossed. Let A be the event
that denotes HEAD on both the throws, B be the event that denotes HEAD on
the first throw, and C be the event that denotes HEAD on the second throw.
Which of the following statements is/are TRUE? \par}
\par\vspace{2pt}\noindent
\textbf{(A)}~A and B are independent.\par\vspace{2pt}\noindent
\textbf{(B)}~A and C are independent.\par\vspace{2pt}\noindent
\textbf{(C)}~B and C are independent.\par\vspace{2pt}\noindent
\textbf{(D)}~Prob(B|C) = Prob(B) \newline CS
\par\vspace{6pt}\hrule height 0.25pt\vspace{6pt}
\noindent\textbf{Q.54} \hfill {\footnotesize\textsf{[MSQ]}}
\par\vspace{2pt}
{\bfseries Consider functions Function 1 and Function 2 expressed in pseudocode as
follows:
Function 1
while n > 1 do
for i = 1 to n do
x = x + 1;
end for
n = ⌊n/2⌋;
end while
Function 2
for i = 1 to 100 ∗n do
x = x + 1;
end for
Let f1(n) and f2(n) denote the number of times the statement ``x = x + 1'' is
executed in Function 1 and Function 2, respectively.
Which of the following statements is/are TRUE? \par}
\par\vspace{2pt}\noindent
\textbf{(A)}~f1(n) $\in$Θ(f2(n))\par\vspace{2pt}\noindent
\textbf{(B)}~f1(n) $\in$o(f2(n))\par\vspace{2pt}\noindent
\textbf{(C)}~f1(n) $\in$$\omega$(f2(n))\par\vspace{2pt}\noindent
\textbf{(D)}~f1(n) $\in$O(n) \newline CS
\par\vspace{6pt}\hrule height 0.25pt\vspace{6pt}
\noindent\textbf{Q.55} \hfill {\footnotesize\textsf{[MSQ]}}
\par\vspace{2pt}
{\bfseries Let G be a simple, finite, undirected graph with vertex set \{v1, . . . , vn\}. Let Δ(G)
denote the maximum degree of G and let N = \{1, 2, . . .\} denote the set of all possible
colors. Color the vertices of G using the following greedy strategy:
for i = 1, . . . , n
color(vi) ←min\{j $\in$N :
no neighbour of vi is colored j\}
Which of the following statements is/are TRUE? \par}
\par\vspace{2pt}\noindent
\textbf{(A)}~This procedure results in a proper vertex coloring of G.\par\vspace{2pt}\noindent
\textbf{(B)}~The number of colors used is at most Δ(G) + 1.\par\vspace{2pt}\noindent
\textbf{(C)}~The number of colors used is at most Δ(G).\par\vspace{2pt}\noindent
\textbf{(D)}~The number of colors used is equal to the chromatic number of G. \newline CS
\par\vspace{6pt}\hrule height 0.25pt\vspace{6pt}
\noindent\textbf{Q.56} \hfill {\footnotesize\textsf{[NAT]}}
\par\vspace{2pt}
{\bfseries Let U = \{1, 2, 3\}. Let 2U denote the powerset of U. Consider an undirected graph
G whose vertex set is 2U. For any A, B $\in$2U, (A, B) is an edge in G if and only
if (i) A̸ = B, and (ii) either A ⊊B or B ⊊A. For any vertex A in G, the set of
all possible orderings in which the vertices of G can be visited in a Breadth First
Search (BFS) starting from A is denoted by B(A).
If ∅denotes the empty set, then the cardinality of B(∅) is
.
CS \par}
\par\vspace{3pt}\noindent\textbf{Answer:} \rule{3.2cm}{0.4pt}
\par\vspace{6pt}\hrule height 0.25pt\vspace{6pt}
\noindent\textbf{Q.57} \hfill {\footnotesize\textsf{[NAT]}}
\par\vspace{2pt}
{\bfseries Consider the following two-dimensional array D in the C programming language,
which is stored in row-major order:
int D[128][128];
Demand paging is used for allocating memory and each physical page frame holds
512 elements of the array D. The Least Recently Used (LRU) page-replacement
policy is used by the operating system. A total of 30 physical page frames are
allocated to a process which executes the following code snippet:
for (int i = 0; i < 128; i++)
for (int j = 0; j < 128; j++)
D[j][i] *= 10;
The number of page faults generated during the execution of this code snippet is
.
CS \par}
\par\vspace{3pt}\noindent\textbf{Answer:} \rule{3.2cm}{0.4pt}
\par\vspace{6pt}\hrule height 0.25pt\vspace{6pt}
\noindent\textbf{Q.58} \hfill {\footnotesize\textsf{[NAT]}}
\par\vspace{2pt}
{\bfseries Consider a computer system with 57-bit virtual addressing using multi-level
tree-structured page tables with L levels for virtual to physical address translation.
The page size is 4 KB (1 KB = 1024 B) and a page table entry at any of the levels
occupies 8 bytes.
The value of L is
.
CS \par}
\par\vspace{3pt}\noindent\textbf{Answer:} \rule{3.2cm}{0.4pt}
\par\vspace{6pt}\hrule height 0.25pt\vspace{6pt}
\noindent\textbf{Q.59} \hfill {\footnotesize\textsf{[NAT]}}
\par\vspace{2pt}
{\bfseries Consider a sequence a of elements a0 = 1, a1 = 5, a2 = 7, a3 = 8, a4 = 9, and
a5 = 2. The following operations are performed on a stack S and a queue Q, both
of which are initially empty.
I: push the elements of a from a0 to a5 in that order into S.
II: enqueue the elements of a from a0 to a5 in that order into Q.
III: pop an element from S.
IV: dequeue an element from Q.
V: pop an element from S.
VI: dequeue an element from Q.
VII: dequeue an element from Q and push the same element into S.
VIII: Repeat operation VII three times.
IX: pop an element from S.
X: pop an element from S.
The top element of S after executing the above operations is
.
CS \par}
\par\vspace{3pt}\noindent\textbf{Answer:} \rule{3.2cm}{0.4pt}
\par\vspace{6pt}\hrule height 0.25pt\vspace{6pt}
\noindent\textbf{Q.60} \hfill {\footnotesize\textsf{[NAT]}}
\par\vspace{2pt}
{\bfseries Consider the syntax directed translation given by the following grammar and se-
mantic rules. Here N, I, F and B are non-terminals. N is the starting non-terminal,
and \#, 0 and 1 are lexical tokens corresponding to input letters ``\#'', ``0'' and ``1'',
respectively. X.val denotes the synthesized attribute (a numeric value) associated
with a non-terminal X. I1 and F1 denote occurrences of I and F on the right hand
side of a production, respectively. For the tokens 0 and 1, 0.val = 0 and 1.val = 1.
N
$\to$
I \# F
N.val = I.val + F.val
I
$\to$
I1 B
I.val = (2 I1.val) + B.val
I
$\to$
B
I.val = B.val
F
$\to$
B F1
F.val = 1
2(B.val + F1.val)
F
$\to$
B
F.val = 1
2B.val
B
$\to$
0
B.val = 0.val
B
$\to$
1
B.val = 1.val
The value computed by the translation scheme for the input string
10\#011
is
. (Rounded offto three decimal places)
CS \par}
\par\vspace{3pt}\noindent\textbf{Answer:} \rule{3.2cm}{0.4pt}
\par\vspace{6pt}\hrule height 0.25pt\vspace{6pt}
\noindent\textbf{Q.61} \hfill {\footnotesize\textsf{[NAT]}}
\par\vspace{2pt}
{\bfseries Consider the following table named Student in a relational database. The primary
key of this table is rollNum.
Student
rollNum
name
gender
marks
1
Naman
M
62
2
Aliya
F
70
3
Aliya
F
80
4
James
M
82
5
Swati
F
65
The SQL query below is executed on this database.
SELECT *
FROM Student
WHERE gender = `F' AND
marks > 65;
The number of rows returned by the query is
.
CS \par}
\par\vspace{3pt}\noindent\textbf{Answer:} \rule{3.2cm}{0.4pt}
\par\vspace{6pt}\hrule height 0.25pt\vspace{6pt}
\noindent\textbf{Q.63} \hfill {\footnotesize\textsf{[NAT]}}
\par\vspace{2pt}
{\bfseries CS
\par\smallskip
Consider the language L over the alphabet \{0, 1\}, given below:
L = \{w $\in$\{0, 1\}∗| w does not contain three or more consecutive 1's\}.
The minimum number of states in a Deterministic Finite-State Automaton (DFA)
for L is
.
CS \par}
\par\vspace{3pt}\noindent\textbf{Answer:} \rule{3.2cm}{0.4pt}
\par\vspace{6pt}\hrule height 0.25pt\vspace{6pt}
\noindent\textbf{Q.64} \hfill {\footnotesize\textsf{[NAT]}}
\par\vspace{2pt}
{\bfseries An 8-way set associative cache of size 64 KB (1 KB = 1024 bytes) is used in a
system with 32-bit address. The address is sub-divided into TAG, INDEX, and
BLOCK OFFSET.
The number of bits in the TAG is
.
CS \par}
\par\vspace{3pt}\noindent\textbf{Answer:} \rule{3.2cm}{0.4pt}
\par\vspace{6pt}\hrule height 0.25pt\vspace{6pt}
\noindent\textbf{Q.65} \hfill {\footnotesize\textsf{[NAT]}}
\par\vspace{2pt}
{\bfseries The forwarding table of a router is shown below.
Subnet Number
Subnet Mask
Interface ID
200.150.0.0
255.255.0.0
1
200.150.64.0
255.255.224.0
2
200.150.68.0
255.255.255.0
3
200.150.68.64
255.255.255.224
4
Default
0
A packet addressed to a destination address 200.150.68.118 arrives at the router.
It will be forwarded to the interface with ID
.
CS \par}
\par\vspace{3pt}\noindent\textbf{Answer:} \rule{3.2cm}{0.4pt}
\par\vspace{6pt}\hrule height 0.25pt\vspace{6pt}
\end{multicols}
\vspace{8pt}
\begin{center}
\rule{0.35\linewidth}{0.5pt}\\[2pt]
\textbf{\small --- END OF QUESTION PAPER ---}\\[-8pt]
\rule{0.35\linewidth}{0.5pt}
\end{center}
\vspace{10pt}
\subsection*{OFFICIAL ANSWER KEY --- GATE 2023 (CS)}
\vspace{-4pt}
\noindent{\footnotesize\textit{Official final answer keys published by the GATE Organising Institute. For NAT questions, allowable numerical ranges are displayed where applicable. MTA = Marks To All.}}
\vspace{4pt}
\begin{center}
\small
\resizebox{\linewidth}{!}{%
\begin{tabular}{|c|c|l||c|c|l||c|c|l||c|c|l|}
\hline
\textbf{Q} & \textbf{Type} & \textbf{Key} & \textbf{Q} & \textbf{Type} & \textbf{Key} & \textbf{Q} & \textbf{Type} & \textbf{Key} & \textbf{Q} & \textbf{Type} & \textbf{Key} \\
\hline\hline
1 & MCQ & B & 17 & MCQ & B & 34 & NAT & 10.2 & 53 & MSQ & C, D \\ \hline
2 & MCQ & MTA & 18 & MCQ & B & 35 & NAT & 7 & 54 & MSQ & A, D \\ \hline
3 & MCQ & A & 20 & MCQ & C & 36 & MCQ & A & 55 & MSQ & A, B \\ \hline
4 & MCQ & A & 22 & MSQ & B, C, D & 38 & MCQ & C & 56 & NAT & 5040 \\ \hline
5 & MCQ & C & 23 & MSQ & C, D & 39 & MCQ & B & 57 & NAT & 4096 \\ \hline
6 & MCQ & D & 24 & MSQ & A, C, D & 41 & MCQ & B & 58 & NAT & 5 \\ \hline
7 & MCQ & C & 25 & MSQ & A, C & 45 & MCQ & C & 59 & NAT & 8 \\ \hline
8 & MCQ & D & 26 & MSQ & B, C & 46 & MCQ & B & 60 & NAT & 2.374 to 2.376 \\ \hline
9 & MCQ & B & 27 & MSQ & A, C, D or C, D & 47 & MCQ & C & 61 & NAT & 2 \\ \hline
10 & MCQ & A & 28 & MSQ & B, D & 48 & MCQ & D & 63 & NAT & 4 \\ \hline
12 & MCQ & B & 29 & MSQ & A, C & 49 & MSQ & B, C, D & 64 & NAT & 19 \\ \hline
13 & MCQ & D & 31 & NAT & 0 & 50 & MSQ & B, C, D & 65 & NAT & 3 \\ \hline
15 & MCQ & A & 32 & NAT & 110 & 51 & MSQ & A, D &  & &  \\ \hline
16 & MCQ & A & 33 & NAT & 2040 & 52 & MSQ & C, D &  & &  \\ \hline
\end{tabular}%
}
\end{center}
\end{document}